% sec-07-ai-peer-review.tex - Section 7, AI Peer Review.  A main section.
% UPDATE stage.  Figures 21, 22, 23 and 24 drawn; Tables 21 to 24 populated.
% Primary source is the AI peer review archive under new-trial-system/inputs;
% the role assignment comes from funding/RFA-RM-27-001-v2.
%
% WHAT THE UPDATE STAGE CHANGED HERE.  The final stage argued the case on
% latency and reviewer count, which is true but not actionable: a funder cannot
% act on "faster".  Two subsections are added that state the case in quantities
% a reader can check and act on within the current fiscal year.
%
%   "What the delay costs" gives the four costs the prior regime imposes, each
%   with a number: US pancreatic cancer deaths inside one review round, the
%   display items a journal article can carry against the fifty this paper
%   carries, the article processing charge against the inference spend, and the
%   direction the money runs between an author and a journal.
%
%   "What three days buys" gives the three the new regime removes: this paper
%   built inside three days, revision at the 1 to 2 day scale evidenced by this
%   very stage, and the author's position that on the record from late 2025 to
%   the present, human peer review is not required for this class of work.
%
% Table 21 is re-cut to those axes and Figure 22's grid is re-cut to match it
% row for row, so the figure and the table cannot drift apart.  The manufacturer
% roster (Table 22), the concurrency figure (Figure 23) and the disagreement
% tree (Figure 24) are carried forward unchanged, by instruction: the Anthropic,
% OpenAI and Google role assignment is the spine of the section.
\section{AI Peer Review}
\label{sec:peerreview}

\subsection{The study the method comes from}

The AI peer review method used throughout this repository was established in a
study deposited November 30, 2025 \cite{aiprgliomatc}. That study was about
clinical trial patient matching in glioblastoma, not pancreatic cancer. The
disease is not what transfers; the review method is, and it transfers without
modification to the pancreatic cancer work reported in \S\ref{sec:ind} through
\S\ref{sec:funding}. The study states its own central claim in one sentence: ``A
primary benefit of AI peer review is the ability to appropriately conduct
iterative analyses and corrections throughout the entire manuscript process.
This responsible method improves on over-worked and fatigued human
recommendations provided after manuscript completion'' \cite{aiprgliomatc}.
Everything else in this section is that claim's measurement. The prior system's
timings are the study's own: the late 1990s and 2000s brought electronic mail
and web portals ``to reduce physical mail delays by up to weeks, for a best-case
review cycle of 7-8 weeks'', and an average journal publication today ``can
still take several months for processing, including 1-2 rounds of review'', with
faster online journals at about one month \cite{aiprgliomatc}. Figure 21 lays
those timings against the new system's on one clock.

% ---------------------------------------------------------------------------
% FIGURE 21.  mermaid-type, sequenceDiagram.  Two lanes carrying six rows each
% over an identical 3.10 cm span, separated by one charcoal hairline, so the
% only difference the figure can express is the axis unit.
% ---------------------------------------------------------------------------
\begin{appfloat}[!tb]
\begin{appfig}[mermaid-type / sequenceDiagram]
\node[mmactor,text width=25mm] (au) at (0.9,0) {Author};
\node[mmactor,text width=25mm] (ed) at (5.0,0) {Editor, prior system};
\node[mmactor,text width=25mm] (hr) at (9.2,0) {Human reviewers, 2 to 3};
\node[mmactor,text width=25mm] (ai) at (13.5,0) {Model reviewers, 3 makers};
\foreach \x in {0.9,5.0,9.2,13.5}{\draw[mmlife] (\x,-0.55) -- (\x,-9.55);}
\begin{scope}[on background layer]
\fill[trialmistl] (-0.5,-4.95) rectangle (14.3,-0.75);
\end{scope}
\node[mmgray,text width=52mm,minimum height=4mm] at (5.05,-0.95) {Prior system: best case 7 to 8 weeks, typical several months};
\seqrow{-1.55}{0.9}{5.0}{mmmsg}{1 submit completed manuscript}
\seqrow{-2.15}{5.0}{9.2}{mmmsg}{2 invite reviewers, wait for acceptance}
\node[mmact,minimum height=1.20cm,fill=trialmist] at (9.2,-2.75) {};
\seqrow{-3.35}{9.2}{5.0}{mmret}{3 reports after weeks, sometimes months}
\seqrow{-3.95}{5.0}{0.9}{mmret}{4 decision plus reports, round 1}
\seqrow{-4.35}{0.9}{5.0}{mmmsg}{5 revision, round 2 begins}
\seqrow{-4.75}{5.0}{0.9}{mmret}{6 second decision, work already finished}
\draw[trialchar,line width=0.5pt] (-0.5,-5.10) -- (14.3,-5.10);
\node[mmsoft,text width=58mm,minimum height=4mm] at (7.20,-5.45) {New system: review during development, hour scale};
\seqrow{-6.05}{0.9}{13.5}{mmmsg}{7 draft plus dataset plus code, mid project}
\seqrow{-6.65}{13.5}{0.9}{mmret}{8 three independent reports, same day}
\seqrow{-7.25}{0.9}{13.5}{mmmsg}{9 corrections applied, re review requested}
\seqrow{-7.85}{13.5}{0.9}{mmret}{10 consensus and recorded disagreement}
\seqrow{-8.55}{0.9}{13.5}{mmmsg}{11 final pass before deposit}
\seqrow{-9.15}{13.5}{0.9}{mmret}{12 release recommendation, human decides}
\node[mmact,minimum height=3.10cm] at (13.5,-7.60) {};
\draw[decorate,decoration={brace,amplitude=4pt},trialchar,line width=0.5pt]
  (-0.75,-1.55) -- (-0.75,-4.75) node[midway,left=3pt,font=\tiny\sffamily,align=right] {weeks to\\months};
\draw[decorate,decoration={brace,amplitude=4pt},trialchar,line width=0.5pt]
  (-0.75,-6.05) -- (-0.75,-9.15) node[midway,left=3pt,font=\tiny\sffamily,align=right] {hours};
\pnote{-0.75}{-10.05}{Both lanes carry six rows over an identical 3.10 cm span. The row counts are
equal by construction, so the only\\thing the two braces can differ on is the
unit on the axis, which is the figure's whole argument.}
\end{appfig}
\vspace{-0.6cm}
\figcaption{Figure 21. One manuscript on two clocks: the prior system's serial
human\\rounds over months, and the new system's parallel model rounds within
hours.}
\end{appfloat}

\subsection{What the delay costs}

Speed is the weakest form of this argument, because a reader cannot act on an
adjective. The four costs the prior regime imposes can be stated as quantities
instead, and each is being paid now rather than in principle.

\textbf{The first cost is measured in patients.} Pancreatic cancer is the third
leading cause of cancer death in the United States, and the 2025 statistics put
US pancreatic cancer deaths at about 51,980 for the year, against a total
five-year survival below 13 percent \cite{siegel2025statistics}. That is about
142 deaths a day and about 1,000 a week. A best-case review round in the prior
regime is 7 to 8 weeks, or about 52 days, so roughly 7,400 Americans die of this
disease inside one best-case round; a typical processing time of several months
covers 17,000 to 26,000. Peer review does not cause those deaths. What the
arithmetic establishes is narrower: the latency of a review regime is a clinical
quantity for this disease, not an administrative one. A dossier, a protocol
amendment or a funding decision that clears in a day rather than in a quarter
reaches a different set of living patients, and that set is countable.

\textbf{The second cost is measured in what a paper is allowed to show.} The
prior regime holds a research article to six or eight figures and tables
combined, for a reason that was once good: every figure cost a person hours to
draw and a reviewer minutes to check. This paper carries twenty-five figures and
twenty-five tables, because each figure here is a specification file rather than
an image: twenty-five files, about 158 KB in total and about 6.5 KB each, every
one carrying a perspective statement, valid source in its platform's own syntax,
an absolute coordinate construction table and an edge-routing argument, as
\S\ref{sec:methods} sets out. A reviewer who doubts Figure 14 reads its
coordinates rather than measuring a rendered image. Substantial, quickly
generated visual evidence is what the prior regime is least equipped to accept,
and it is what a funder asks for first.

\textbf{The third cost is measured in dollars and in months of lost research.}
List article processing charges commonly run from about \$2,000 to \$6,000 per
article and exceed \$10,000 at several flagship titles, and the researcher pays
them. The comparison figure from the source study is exact: its evaluation
standard scored a perfect project completion against ``n\textsubscript{p} = 5
prompts within n\textsubscript{d} = 12 days while only spending
n\textsubscript{c} = \$35'' \cite{aiprgliomatc}. Thirty-five dollars of
inference carried a fourteen-day study through a dataset, two notebooks, six
regulatory documents and a working API repository. The months matter more than
the dollars. An author placing three papers a year in the prior regime spends
twenty-one to twenty-four weeks a year waiting on rounds, and the source study
names what those weeks are: the new regime
``eliminates publication lull and cognitive load of uncertainty, as the
researcher simply transitions to the next AI peer reviewed project''
\cite{aiprgliomatc}.

\textbf{The fourth cost is the direction the money runs.} Under the author-pays
model the researcher transfers the charge to the journal, and what the journal
supplies in return is a page on its own site carrying the researcher's work,
with the indexing, the traffic and the promotion of that work accruing to the
journal's own title. The source study states the alternative: authors ``can now improve their speed of development and gain
reputation in papers by simply providing prompts, LLMs, AI peer review
recommendations, and corrections they made - without the need for paid review or
journals'', and ``interested parties can then view the author's peer review
findings alongside the paper'' \cite{aiprgliomatc}. Every artifact in this
repository is deposited to a DOI and to a public repository at no charge to the
author, under CC BY, with the review findings deposited beside the work rather
than withheld from it. Table 21 puts the four costs beside what replaces each of
them, and Figure 22 computes the ratio.

\begin{apptable}
\begin{tabularx}{\textwidth}{>{\raggedright\arraybackslash}p{3.70cm} >{\raggedright\arraybackslash}p{3.65cm} >{\raggedright\arraybackslash}p{3.70cm} >{\raggedright\arraybackslash}X}
\headrow \thead{Axis} & \thead{Prior system} & \thead{New system} & \thead{Ratio or delta}\\
\midrule
Latency, one round & 7 to 8 weeks best case & Same day & About 50 to 1 \\
Typical journal processing & Several months, 1 to 2 rounds & Hours, rounds unlimited & About 300 to 1 \\
US pancreatic cancer deaths inside one round & About 7,400 over 52 days & About 6 over 1 hour & The cost of the wait \\
Display items one paper carries & Commonly 6 to 8 per article & 25 figures and 25 tables & About 3 to 1 \\
Author cash cost per paper & \$2,000 to \$6,000 list charge, over \$10,000 at flagship titles & \$35 of inference across a 14-day study & About 57 to 1 and up \\
Direction the money runs & Author pays the journal to host the author's work & Author pays no journal, deposits under CC BY & Not comparable \\
Reviewers per round & 2 to 3 humans & 3 model manufacturers & 1 to 1 by count \\
Artifacts absorbed per year & 2 to 6 per author & Over 30 deposited in 2026 & About 6 to 1 \\
Corrections before release & None, review follows release & Every round, before deposit & The decisive axis \\
\end{tabularx}
\end{apptable}
\vspace{-0.6cm}
\tabcap{Table 21. Eight measurable axes of one review round under both regimes,\\
and the ninth axis on which the two regimes are not comparable at all.}

Figure 22 carries seven of those axes as a grid, with the ratio computed in each
cell and the decisive row separated by its own rule. A regime whose review
arrives after release cannot apply a correction to the released object, whatever
its cost or its latency, and that is why the last row is not an economic axis at
all.

% ---------------------------------------------------------------------------
% FIGURE 22.  d2-type, grid.  UPDATE STAGE RE-CUT.  Seven data rows carrying the
% axes of Table 21, one row height, four fixed column widths, no edges, cut
% twice: once below the header and once above the decisive row.  The rows are
% held in the same order as the table so the two cannot drift apart.
% ---------------------------------------------------------------------------
\begin{appfloat}[!tb]
\begin{appfig}[d2-type / grid]
\def\rh{0.64}
\node[d2cellh,minimum width=4.4cm,minimum height=\rh cm,anchor=north west] at (0,0) {Axis};
\node[d2cellh,minimum width=4.0cm,minimum height=\rh cm,anchor=north west] at (4.4,0) {Prior system, human review};
\node[d2cellh,minimum width=4.0cm,minimum height=\rh cm,anchor=north west] at (8.4,0) {New system, AI review};
\node[d2cellh,minimum width=2.2cm,minimum height=\rh cm,anchor=north west] at (12.4,0) {Ratio};
\foreach \i/\a/\p/\n/\r/\rs in {%
  1/{Latency, one round}/{7 to 8 weeks best case}/{same day}/{about 50 to 1}/{d2cellk},
  2/{Journal processing}/{several months, 1 to 2 rounds}/{hours, rounds unlimited}/{about 300 to 1}/{d2cellk},
  3/{US PDAC deaths per round}/{about 7,400 over 52 days}/{about 6 over 1 hour}/{the cost of the wait}/{d2cellk},
  4/{Display items per paper}/{commonly 6 to 8 per article}/{25 figures and 25 tables}/{about 3 to 1}/{d2cellk},
  5/{Author cash per paper}/{\$2,000 to \$6,000 list charge}/{\$35 of inference, 14-day study}/{about 57 to 1}/{d2cellk},
  6/{Direction the money runs}/{author pays the journal}/{author pays no journal}/{not comparable}/{d2cellg},
  7/{Artifacts absorbed per year}/{2 to 6 per author}/{over 30 deposited in 2026}/{about 6 to 1}/{d2cellk}}{
  \node[d2cellg,minimum width=4.4cm,minimum height=\rh cm,anchor=north west] at (0,{-\rh*\i}) {\a};
  \node[d2cell,minimum width=4.0cm,minimum height=\rh cm,anchor=north west] at (4.4,{-\rh*\i}) {\p};
  \node[d2celll,minimum width=4.0cm,minimum height=\rh cm,anchor=north west] at (8.4,{-\rh*\i}) {\n};
  \node[\rs,minimum width=2.2cm,minimum height=\rh cm,anchor=north west] at (12.4,{-\rh*\i}) {\r};}
\node[d2cellk,minimum width=4.4cm,minimum height=\rh cm,anchor=north west] at (0,{-\rh*8}) {Corrections before release};
\node[d2cellg,minimum width=4.0cm,minimum height=\rh cm,anchor=north west] at (4.4,{-\rh*8}) {none, review follows release};
\node[d2celll,minimum width=4.0cm,minimum height=\rh cm,anchor=north west] at (8.4,{-\rh*8}) {every round, before deposit};
\node[d2cellh,minimum width=2.2cm,minimum height=\rh cm,anchor=north west] at (12.4,{-\rh*8}) {decisive};
\draw[trialchar,line width=0.7pt] (0,{-\rh}) -- (14.6,{-\rh});
\draw[trialchar,line width=0.7pt] (0,{-\rh*8}) -- (14.6,{-\rh*8});
\node[d2mid,text width=54mm,anchor=north west] at (0,{-\rh*9-0.20}) {About 1,000 US pancreatic cancer deaths per week of review latency};
\node[d2soft,text width=50mm,anchor=north west] at (5.90,{-\rh*9-0.20}) {\$35 of inference against a \$2,000 to \$10,000 charge};
\node[d2gray,text width=48mm,anchor=north west] at (11.20,{-\rh*9-0.20}) {50 display items here against 6 to 8 allowed};
\pnote{0}{-7.55}{The last row is separated by its own rule because it is not an economic axis: it
is the reason the economics\\matter. A regime whose review follows release cannot
correct the object that was released.}
\end{appfig}
\vspace{-0.6cm}
\figcaption{Figure 22. Seven measurable axes of one review round under both
regimes,\\with the ratio in each cell and the decisive row cut away below them.}
\end{appfloat}

\subsection{What three days buys}

Three things follow from the four costs, and each is stated here as a claim with
a date attached rather than as a prospect.

\textbf{This paper was finished within three days.} The master prompt was
supplied on August 12, 2026 and this document was deposited on August 14, 2026,
with twenty-five figures, twenty-five tables, eleven section files and about
324,000 characters of LaTeX source, across an eight-stage build whose every file
was committed the moment it was written. Nothing in that window waited on
anyone. The point is not that three days beats several months; it is that the
prior regime has no state in which a document of this size and this figure count
exists at all before the first review round has closed.

\textbf{Major revisions run on the 1 to 2 day scale.} This version of the paper
is the evidence for that claim rather than an assertion of it. A second author
prompt on August 13, 2026 required Figure 14 to be redrawn from a use case
diagram to a class diagram, this section to be rewritten around quantities,
Table 21 and Figure 22 to be re-cut to the same axes, and a defect that stopped
the prior stage compiling at all to be found and fixed; the result was deposited
the following day. In the prior regime a change of that size is a revision
round, and a revision round is 7 to 8 weeks at best. The new regime completes
about thirty of them in the time the prior regime completes one, and completes
them before release rather than after it.

\textbf{On this record, human peer review is not required for this class of
work.} That is the author's position and it is stated plainly, because hedging
it would misrepresent how the repository is actually run. The basis is the work
from late 2025 to the present: five separately deposited and separately dated
review instances, listed in Table 24, whose corrections are visible in the
diffs; frontier model results on academic reasoning, scientific knowledge and
multimodal benchmarks that improved across every release in that window
\cite{aiprgliomatc}; and a production rate that the prior regime's throughput
cannot meet at any level of goodwill. The position is bounded, and the three
bounds are set out below under what AI peer review is not: it concerns journal
peer review of the author's own deposited works. It does not touch institutional
review board
approval, an FDA information request, a data and safety monitoring board, or any
other regulatory determination, none of which is a peer review process and none
of which this system proposes to replace.

\subsection{Three manufacturers, one human gate}

The pancreatic cancer funding application makes the role assignment explicit and
binding. Its provenance section states that the ``planned tripartisan workflow
assigns Claude Code the primary production role for code, regulatory
adaptations, protocols, IND materials, diagrams, and repository-scale artifacts.
ChatGPT Codex serves as the first independent peer reviewer, emphasizing
verification, validation, uncertainty quantification, external standards,
trial-literature synthesis, deep research, and PDF review. Google Gemini serves
as the second independent reviewer, emphasizing accelerated code problem
solving, meta-verification of prior stages, and administrative workflow
support'' \cite{rfarm27001v2}. The reviews ``occur at defined mid-project
milestones and after material corrections so that findings can change the work
before release, with the human PI retaining final authority'', and the same
paragraph records that these systems ``are disclosed as drafting and peer-review
aids'' \cite{rfarm27001v2}. Table 22 names the three manufacturers, the role
each holds, and the human principal investigator who retains final authority
over all three; Figure 23 draws where the concurrency sits and where it stops.
Three manufacturers read one frozen artifact and its recorded hash, so no branch
can observe another's output and a disagreement is evidence rather than an
artifact of ordering.

\begin{apptable}
\begin{tabularx}{\textwidth}{>{\raggedright\arraybackslash}p{2.55cm} >{\raggedright\arraybackslash}p{3.05cm} >{\raggedright\arraybackslash}X}
\headrow \thead{Manufacturer} & \thead{Role} & \thead{What that role covers in this work}\\
\midrule
Anthropic & Primary production & Code, regulatory adaptations, protocols, IND materials, diagrams, repository-scale artifacts \\
OpenAI & First independent reviewer & Verification, validation, uncertainty quantification, external standards, trial-literature synthesis, deep research, PDF review \\
Google & Second independent reviewer & Accelerated code problem solving, meta-verification of prior stages, administrative workflow support \\
Human PI & Final authority & Source verification, regulatory accuracy, scientific interpretation, participant protection, final submission \\
\end{tabularx}
\end{apptable}
\vspace{-0.6cm}
\tabcap{Table 22. The three model manufacturers and the human principal\\
investigator, and the role each of them holds in the reviewed work.}

% ---------------------------------------------------------------------------
% FIGURE 23.  plantuml-type, activity with fork and join.  Three branches at a
% 5.20 cm pitch, 18 mm of clear canvas either side of each, one human gate.
% ---------------------------------------------------------------------------
\begin{appfloat}[!tb]
\begin{appfig}[plantuml-type / activity with fork and join]
\node[umlinit] (i0) at (7.50,0) {};
\node[umlbox,text width=46mm] (a1) at (7.50,-0.85) {Artifact reaches a defined mid project milestone};
\node[umlbox,text width=46mm] (a2) at (7.50,-2.05) {Freeze the artifact and its inputs, record the hash};
\node[umlbar,minimum width=12.8cm] (fk) at (7.50,-2.95) {};
\node[umlkey,text width=34mm] (b1) at (2.30,-4.45) {Anthropic, production role: code, regulatory text, diagrams};
\node[umlctrl,text width=34mm] (b2) at (7.50,-4.45) {OpenAI, first reviewer: verification, validation, uncertainty, standards};
\node[umlctrl,text width=34mm] (b3) at (12.70,-4.45) {Google, second reviewer: code problem solving and meta verification};
\node[umlbar,minimum width=12.8cm] (jn) at (7.50,-6.05) {};
\node[umlbox,text width=42mm] (a3) at (7.50,-6.85) {Join, three reports over one frozen artifact};
\node[mmdec,aspect=2.2,text width=24mm] (d1) at (7.50,-7.95) {Reports agree?};
\node[umlbox,text width=36mm] (cn) at (3.60,-9.05) {Apply the consensus corrections};
\node[umlstategray,text width=36mm] (dg) at (11.40,-9.05) {Record the disagreement verbatim, both positions};
\node[umlstategray,text width=36mm] (rt) at (11.40,-10.05) {Route to the Figure 24 resolution tree};
\node[umlkey,text width=46mm,line width=1pt] (pi) at (7.50,-11.15) {Human principal investigator reviews, approves or rejects};
\node[mmdec,aspect=2.2,text width=22mm] (d2) at (7.50,-12.35) {Approved?};
\node[umlbox,text width=32mm] (rl) at (4.20,-13.45) {Release, deposit, provenance recorded};
\node[umlfinal] (f1) at (4.20,-14.35) {};
\fill[trialburg] (f1.center) circle (1.1mm);
\node[umlbox,text width=32mm] (rj) at (11.40,-13.45) {Return to production with the rejection reason};
\node[umlfinal] (f2) at (11.40,-14.35) {};
\draw[trialchar,line width=0.6pt] ([shift={(-0.13,-0.13)}]f2.center) -- ([shift={(0.13,0.13)}]f2.center);
\draw[umlarrow] (i0) -- (a1);
\draw[umlarrow] (a1) -- (a2);
\draw[umlarrow] (a2) -- (fk);
\foreach \n in {b1,b2,b3}{\draw[umlarrow] (\n |- fk.south) -- (\n.north);
  \draw[umlarrow] (\n.south) -- (\n |- jn.north);}
\draw[umlarrow] (jn) -- (a3);
\draw[umlarrow] (a3) -- (d1);
\draw[umlarrow] (d1.west) -- ++(-0.55,0) -- (3.60,-7.95) -- node[umlguard,pos=0.6] {yes} (cn.north);
\draw[umlarrow] (d1.east) -- ++(0.55,0) -- (11.40,-7.95) -- node[umlguard,pos=0.6] {no} (dg.north);
\draw[umlarrow] (dg) -- (rt);
\draw[umlarrow] (cn.south) -- (3.60,-10.55) -- (pi.north west);
\draw[umlarrow] (rt.south) -- (11.40,-10.55) -- (pi.north east);
\draw[umlarrow] (pi) -- (d2);
\draw[umlarrow] (d2.west) -- ++(-0.55,0) -- (4.20,-12.35) -- node[umlguard,pos=0.6] {yes} (rl.north);
\draw[umlarrow] (d2.east) -- ++(0.55,0) -- (11.40,-12.35) -- node[umlguard,pos=0.6] {no} (rj.north);
\draw[umlarrow] (rl) -- (f1);
\draw[umldash] (rj) -- (f2);
\pnote{0}{-15.15}{The three branches sit at a 5.20 cm pitch against a 34 mm text width, leaving
18 mm of clear canvas either side, so\\the fork and join bars read as single
horizontal strokes rather than as a hatched band.}
\end{appfig}
\vspace{-0.6cm}
\figcaption{Figure 23. One artifact reviewed concurrently by three model
manufacturers,\\joined into one recorded disagreement set, and cleared only by the
human PI.}
\end{appfloat}

\subsection{What the review actually found}

The glioblastoma study is worth reporting in its numbers, because the numbers
are what make the method transferable rather than aspirational. A ten-file
dataset and a clinical BERT notebook reached ``a final notebook at a low test
set performance of 67.3\%''. A triple review by three models from three
manufacturers ``all yielded the primary recommendation to upgrade to a tabular
model, which was more suitable for the csv dataset''; one of them identified an
open-source TabPFN-v2 classifier, and the recommendations were implemented to
yield ``a test set accuracy improvement to 94.0\% with few-shot learning along
with peer review corrections'' \cite{aiprgliomatc}. An evaluation standard
generated for the purpose scored the workflow at 87.5 percent overall quality
and 0.831 efficiency. Eleven attempts to produce a working FastAPI repository
failed before a two-model prompting strategy succeeded. Three independent
regulatory analysts scored the six generated FDA, ICH and ISO documents at 8.0,
9.1 and 9.1 out of 10, and the study records the reason each gave rather than
averaging them \cite{aiprgliomatc}. Table 23 collects the measures.

\begin{apptable}
\begin{tabularx}{\textwidth}{>{\raggedright\arraybackslash}p{5.35cm} >{\raggedright\arraybackslash}p{2.55cm} >{\raggedright\arraybackslash}X}
\headrow \thead{Measure} & \thead{Value} & \thead{What changed it}\\
\midrule
Test-set accuracy before review & 67.3 percent & Initial clinical BERT pipeline \\
Test-set accuracy after review & 94.0 percent & Tabular model plus few-shot learning \\
Overall workflow quality & 87.5 percent & Evaluation standard, eight criteria \\
Efficiency score & 0.831 & Human, prompt, time and cost metrics \\
Regulatory analyst scores & 8.0, 9.1, 9.1 of 10 & Three manufacturers, recorded separately \\
Failed API attempts before success & 11 & Resolved by a two-model prompting strategy \\
Study duration, one author & 14 days & Dataset through FastAPI repository \\
\end{tabularx}
\end{apptable}
\vspace{-0.6cm}
\tabcap{Table 23. What the AI peer review changed, measured, in the glioblastoma\\study
whose method the pancreatic cancer work in this paper inherits.}

Figure 24 is the reason those three regulatory scores are reported separately
rather than as a mean of 8.7. None of its five terminal dispositions averages
the reports, only one leaves a disagreement without an action, and two stop the
artifact from progressing at all.

% ---------------------------------------------------------------------------
% FIGURE 24.  graphviz-type, decision tree.  Four ranks at a uniform 2.05 cm
% separation, deliberately asymmetric, with a 1.50 cm empty corridor between the
% two rank-3 subtrees.
% ---------------------------------------------------------------------------
\begin{appfloat}[!tb]
\begin{appfig}[graphviz-type / decision tree]
\node[gvboxk,text width=52mm,line width=0.9pt] (rt) at (7.50,0) {Three reports over one frozen artifact};
\node[mmdec,aspect=2.0,text width=22mm] (d1) at (7.50,-2.05) {Do all three agree?};
\node[gvboxs,text width=30mm] (t1) at (1.85,-4.10) {Apply the consensus};
\node[mmdec,aspect=2.0,text width=24mm] (d2) at (9.35,-4.10) {Factual or judgmental?};
\node[mmdec,aspect=2.0,text width=24mm] (d3) at (6.05,-6.15) {Checkable against a cited source?};
\node[mmdec,aspect=2.0,text width=24mm] (d4) at (12.65,-6.15) {Changes a safety claim?};
\node[gvboxs,text width=30mm] (t2) at (3.55,-8.20) {Resolve to the source, record which model erred};
\node[gvboxm,text width=30mm] (t5) at (7.35,-8.20) {Hold the artifact, re-run the round};
\node[gvboxm,text width=30mm] (t3) at (11.15,-8.20) {Escalate to the human PI, both positions verbatim};
\node[gvboxs,text width=30mm] (t4) at (14.15,-8.20) {Record and proceed, note in the limitations};
\draw[gvedge] (rt) -- (d1);
\draw[gvedge] (d1) -- node[umlguard,pos=0.5] {yes} (t1);
\draw[gvedge] (d1) -- node[umlguard,pos=0.5] {no} (d2);
\draw[gvedge] (d2) -- node[umlguard,pos=0.5] {factual} (d3);
\draw[gvedge] (d2) -- node[umlguard,pos=0.5] {judgmental} (d4);
\draw[gvedge] (d3) -- node[umlguard,pos=0.5] {yes} (t2);
\draw[gvedgeb] (d3) -- node[umlguard,pos=0.5] {no} (t5);
\draw[gvedgeb] (d4) -- node[umlguard,pos=0.5] {yes} (t3);
\draw[gvedge] (d4) -- node[umlguard,pos=0.5] {no} (t4);
\node[gvboxg,text width=44mm] at (0.35,-6.15) {No disposition averages the reports. Two of five stop the artifact.};
\pnote{0}{-9.45}{The tree is deliberately asymmetric: the left branch terminates in one step
because agreement is the common case,\\and the right branch in three because
disagreement is the case the figure exists to document.}
\end{appfig}
\vspace{-0.6cm}
\figcaption{Figure 24. One disagreement from three model reports to five terminal\\
dispositions, of which none is an average and only one is silent.}
\end{appfloat}

\subsection{Why only AI reviewers can carry this load}

The study's conclusion is not that AI review is convenient. It is that the
volume, complexity and frequency of AI-generated research have already passed
what human review can absorb: this next evolution of peer review marks ``the
pivot away from slow and irresponsible use of human peer review for
LLM-generated data, and back to fast automation of patient health applications
at scale'' \cite{aiprgliomatc}. The same study records that trust for the
process ``was established over a series of human-AI works reflected in LLM code
and non-code benchmarking, applications in literature, and prior author works
using AI in review roles coordinated across judging, external validation,
cross-verification, meta-verification, and tests according to FDA/ICH guidance''
\cite{aiprgliomatc, e2eefficiency2025, gbmsynergy2025}. Table 24 lists five
instances of that process across the author's deposited works, and the specific
change each produced. The count itself is the argument: in 2024 the author
deposited eight works, in 2025 twelve, and in 2026 through August more than
twenty, including an IND, two clinical trial protocols, five bill versions, ten
funding applications, two NIH application versions and a capitalization plan
\cite{repo2026}. A regime that took seven to eight weeks per round in its best
case was already the binding constraint at eight works a year; at thirty it is
not a constraint but an exclusion, and the work simply does not enter it. That
is the sense in which the two systems are incompatible rather than merely
different. Human review is not worse; it is sized for a production regime that
no longer exists on the author's side of the transaction, and no amount of
goodwill inside the prior regime changes its throughput.

\begin{apptable}
\begin{tabularx}{\textwidth}{>{\raggedright\arraybackslash}p{4.95cm} >{\raggedright\arraybackslash}p{1.85cm} >{\raggedright\arraybackslash}p{2.55cm} >{\raggedright\arraybackslash}X}
\headrow \thead{Work} & \thead{Date} & \thead{Reviewers} & \thead{What the review changed}\\
\midrule
Glioblastoma patient matching ML & Nov 30, 2025 & 3 manufacturers & Model class, 67.3 to 94.0 percent \\
Code generation competition & Dec 22, 2025 & 16 models & Iterative learning on FDA adverse events \\
NIH application v1.0 & Jul 7, 2026 & 2 reviewers & Structure and mechanism fit \\
NIH application v2.0 & Jul 12, 2026 & 2 reviewers & Budget, milestones, review schedule \\
Capitalization plan & Aug 11, 2026 & 2 reviewers & Column widths, caption balance, figure centering \\
\end{tabularx}
\end{apptable}
\vspace{-0.6cm}
\tabcap{Table 24. Five instances of AI peer review across the author's\\
deposited works, and the specific change each review produced.}

\subsection{What AI peer review is not}

Three limits bound the position taken above, and the case is weaker if they are
left implicit. First, model review is not regulatory review. Nothing in
\S\ref{sec:peerreview} substitutes for an institutional review board, an FDA
information request, or a data and safety monitoring board, and the funding
application says so in its own words: the models ``are not applicants,
investigators, sponsors, regulators, clinicians, or decision-makers''
\cite{rfarm27001v2}. Second, model review is not clinical validation.
Improvements over time on coding, scientific-reasoning and multimodal benchmarks
support current frontier models as independent reviewers, ``but not as clinical
validators'' \cite{rfarm27001v2, aiprgliomatc}. Third, model review does not
establish correctness by agreement. Three models agreeing is three reports, not
a proof, which is why Figure 24's tree never averages and why two of its five
dispositions stop the artifact rather than passing it. The source study bounds
itself the same way, recording that its own evaluation standard was ``only based
on the current project'' and its efficiency metrics ``were partly influenced by
the author's priorities'' \cite{aiprgliomatc}; that is why Table 23 reports
three regulatory scores separately rather than a single figure of merit.

\subsection{The trust that had to be earned first}

AI peer review is only defensible if the reviewers were shown to be competent
before they were relied upon, and that record runs from late 2025 to the
present. In October 2025 an end-to-end oncology trial efficiency study
established that LLM-generated regulatory text could be produced against FDA and
ICH requirements and checked \cite{e2eefficiency2025}; in November 2025 a drug
synergy machine learning study took rapid code prototypes through to project
deliverables \cite{gbmsynergy2025}; later that month the glioblastoma patient
matching study introduced the triple review this section reports
\cite{aiprgliomatc}; and in December 2025 a code generation competition ran
sixteen proprietary and open-source models against FDA adverse event data with
iterative learning, which measures reviewer competence rather than assuming it
\cite{codegencomp2025}. By the time the pancreatic cancer work of
\S\ref{sec:ind} through \S\ref{sec:funding} began in 2026, the question was no
longer whether the models could review; it was which model should hold which
role, and the answer is Table 22.

The chronology matters because it is the difference between an appeal to
authority and an appeal to record. The claim is not that frontier models are
reliable in general, but that these models, in these roles, on this class of
artifact, produced measurable corrections on five separately deposited and
separately dated occasions whose diffs are public. A reviewer who doubts the
claim can read the corrections rather than the argument.

\subsection{What a human still does}

Four things in this system are done by a person and cannot be delegated. The
person decides what to build: no model chose to write an IND, two protocols, a
bill and fourteen funding artifacts. The person supplies the master prompt and
therefore the constraints, including the ones that make the output checkable,
such as one file per commit and no raster figures. The person adjudicates
disagreement: Figure 24's escalate branch terminates at the principal
investigator, and two of its five dispositions stop the artifact. And the person
signs, which the funding application states in the plainest terms available,
that the models ``are not applicants, investigators, sponsors, regulators,
clinicians, or decision-makers'' \cite{rfarm27001v2}.

That division is the reason the system can be audited at all. A regime in which
a model both produced and cleared the work would have no point at which
responsibility attaches, and no funder should accept one. The regime described
here attaches responsibility at exactly one point, and the whole of
\S\ref{sec:peerreview} argues that the point is correctly placed.
